\appendix

\section{Boundary cases of the gas-exchange closure}
\label{app:closure}

Section~\ref{sec:closure} solves the supply--demand closure as a quadratic in
$C_i$ and returns the admissible positive root on each branch. Two limiting
conductance states need a separate treatment. Both have physical answers rather
than clamps, and neither is exercised in any reported run, but stating them
completes the definition of the closure.

The vanishing-conductance case has a physical answer rather than a clamp, and the
one-sided limit and the closed boundary are two different states. As
$g_c \to 0^{+}$ the quadratic degenerates: $a \to 0$ and
$b \to R_d d_1 - \mathrm{num}$, and the finite root of the remaining linear
equation is

\begin{equation}
\lim_{g_c \to 0^{+}} C_i
= \frac{\mathrm{num}\,\Gamma^{*} + R_d d_0}{\mathrm{num} - R_d d_1},
\label{eq:ci_limit}
\end{equation}

provided the denominator is positive, which holds for both branches at the
parameter values used here. Substituting this root into
Eq.~\eqref{eq:branch_form} gives $A_b = R_d$ and hence $A_n \to 0$: supply and
demand meet at the compensation point, where diffusion delivers nothing and
gross assimilation exactly pays respiration. The limit is a finite
$C_i$, not $\Gamma^{*}$. At the production values it evaluates to
$67.94~\mu\mathrm{mol\,mol^{-1}}$ on the Rubisco-limited branch and
$53.35$ on the electron-transport-limited branch, and a numerical sweep in
$g_c$ reproduces both.

Exactly closed stomata are a separate case. At $g_c = 0$ the supply relation
carries no information about $C_i$, the quadratic is empty, and the physical
state is a leaf that admits no CO$_2$ and only respires. The implementation
returns that state directly at the boundary $g_c \leq 0$,

\begin{equation}
C_i = \Gamma^{*}, \qquad A_n = -R_d ,
\label{eq:ci_dark}
\end{equation}

which is the dark limit and replaces a clamp on the conductance.

This assignment is an implementation boundary convention and not the continuous
extension of Eq.~\eqref{eq:ci_limit}. The one-sided limit derived above gives
$A_n \to 0$ as $g_c \to 0^{+}$, whereas the convention returns $A_n = -R_d$ at
$g_c = 0$, so the two differ by $R_d$ at the boundary. We state the difference
rather than smooth it over. The convention is the physically meaningful one for a
leaf with no gas exchange, since gross assimilation requires a CO$_2$ supply that
a shut stoma does not provide, while the limit describes a leaf whose stomata are
almost shut and which is still exchanging enough CO$_2$ to break even. Whether
the discrepancy matters is a question about the trajectory: the conductance model
of Section~\ref{sec:jarvis} has a positive floor on its light response, so
$g_c = 0$ is never reached in any reported run and the branch is not exercised.
Respiration in darkness is carried by the same expression that carries
assimilation in the light.

One numerical caveat belongs with this derivation. Written in the form above the
quadratic loses significance when $a = -g_c d_1$ becomes comparable to the
rounding error of $b$, which at these parameter values occurs below
$g_c \approx 10^{-8}$; a sweep to $g_c = 10^{-11}$ returns a spurious positive
$A_n$. The conductance model of Section~\ref{sec:jarvis} has a positive floor on
its light response, so no reported run enters that range, and we record the
bound rather than rely on it being unreachable.

\section{Parameter provenance ledger}
\label{app:params}

Every numerical constant in the released model carries a provenance record in
\texttt{config/parameter\_provenance.csv}, reproduced in full as Table~S1 of the
Supplementary Material. The tier states how the value
was obtained: T1 measured, T2 derived, T3 inherited, T4 constrained, T5 fitted, T6 unsourced. A T6 entry is a value that
we could not trace to a measurement, a derivation or the source model, and it is
declared as such rather than presented as sourced. The ledger holds 124 entries, of which 4 are measured (T1), 55 are derived (T2), 27 are inherited (T3), 9 are constrained (T4), 15 are fitted (T5) and 14 are unsourced (T6). Every key of the released configuration appears in the ledger, as does every constant hard-coded in the model source, and every ledger value that names a configuration key matches the value the reported runs load; these properties are checked by script rather than by reading. Two rows are attributed to this work rather than to a publication, because no external source exists for them: the growth coefficient obtained by the held-out sweep and the assimilation scale constrained by the reference data.

The complete ledger of all 124 entries, with value, unit, tier, source and
locator for each, is provided as Table~S1 in the Supplementary Material and is
released with the code as \texttt{config/parameter\_provenance.csv}. The curated
physical parameter table in the main text (Table~\ref{tab:params}) is the subset
most relevant to reading the model equations.

\section{Sensitivity analysis: parameters, ranges and sampling scales}
\label{app:sa}

Table~\ref{tab:sa_ranges} is generated from
\texttt{config/sensitivity\_params.csv} and gives the factors varied in the
sensitivity study together with their ranges, sampling scales and the role each
factor plays in the model. All reported Sobol' and Morris indices are conditional
on exactly these ranges: an index is a statement about output variance over a
declared input distribution, and it is not transferable to a different one. A
factor sampled on a logarithmic scale is varied uniformly in $\log_{10}$, which
is the appropriate choice for a rate constant spanning more than one order of
magnitude. Of the 20 factors listed, 16 are active in the reported design and 4
are disabled; the reasons for disabling are given in the table.

\begin{table*}[!htbp]
\centering
\scriptsize
\caption{Sensitivity analysis design, generated from
\texttt{config/sensitivity\_params.csv}. Role distinguishes a model parameter from
an environmental driver, an initial condition and a physical constant. The basis
column records why the range was chosen.}
\label{tab:sa_ranges}
\begin{tabular}{@{}p{46mm}ccccp{54mm}@{}}
\toprule
Factor & Lower & Upper & Scale & Role & Basis \\
\midrule
\multicolumn{6}{@{}l}{\textit{Active in the reported design}} \\
\texttt{growth.lambda\_sb} & 0.002 & 0.01 & linear & model parameter & Calibrated; $\pm$70\% around default 0.00587 (Schulz et al. 2019) \\
\texttt{carbohydrates.respiration\_frequency} & 0.4 & 1.2 & linear & model parameter & Calibrated; $\pm$50\% around default 0.79 \\
\texttt{nutrients.max\_nitrogen\_uptake} & 1 & 15 & log & model parameter & Literature range for herbaceous crops (Barber 1995; Lambers et al. 2008) \\
\texttt{nutrients.max\_phosphorus\_uptake} & 0.05 & 1.5 & log & model parameter & Literature range for herbaceous crops (Barber 1995) \\
\texttt{nutrients.nitrogen\_soil\_content} & 3 & 30 & linear & environment & Low-nutrient to high-nutrient agricultural soil (Brady \& Weil 2010) \\
\texttt{nutrients.phosphorus\_soil\_content} & 0.03 & 0.5 & log & environment & Low to high-P agricultural soil (Brady \& Weil 2010) \\
\texttt{gas\_exchange.leaf\_temperature} & 10 & 35 & linear & environment & Controlled environment range; cold stress to heat stress \\
\texttt{gas\_exchange.relative\_humidity} & 30 & 90 & linear & environment & Arid to humid growing conditions \\
\texttt{gas\_exchange.ambient\_CO2\_concentration} & 280 & 800 & linear & environment & Pre-industrial to RCP 8.5 projected 2100 (Meinshausen et al. 2011) \\
\texttt{gas\_exchange.leaf\_water\_potential} & -1.5 & -0.1 & linear & environment & Severe water stress to turgid (Hsiao 1973; Jones 2014) \\
\texttt{gas\_exchange.max\_rubisco\_carboxylation\_rate} & 20 & 120 & log & model parameter & C3 herbaceous range (Wullschleger 1993; Medlyn et al. 2002) \\
\texttt{gas\_exchange.electron\_transport\_rate} & 50 & 200 & log & model parameter & C3 herbaceous range (Wullschleger 1993) \\
\texttt{photo.max\_PAR} & 300 & 2000 & linear & environment & Shade canopy (300) to full sun (2000 $\mu$mol m⁻² s⁻¹) (McCree 1972) \\
\texttt{carbohydrates.sucrose} & 0.8 & 3 & linear & initial condition & Observed range in spinach seedlings (\% DW; Geiger \& Servaites 1994) \\
\texttt{growth.total\_biomass} & 0.01 & 0.12 & linear & initial condition & Seedling (0.01 g FW) to established juvenile plant (0.12 g FW) \\
\texttt{carbohydrates.k\_vacuole\_in} & 0.005 & 0.4 & linear & model parameter & Cytosol→vacuole exchange rate; covers calibrated per-PP range 0.011-0.230 h-1 (Lohaus 2007) \\
\midrule
\multicolumn{6}{@{}l}{\textit{Disabled}} \\
\texttt{carbohydrates.lambda\_sni} & 0.04 & 0.18 & linear & model parameter & Disabled: mu*$\approx$0 everywhere; excluded to keep d$\le$16 (solver limit) \\
\texttt{photo.feedback\_on\_photosynthesis} & 0.5 & 0.9 & linear & model parameter & Disabled: inactive in Farquhar path (light\_PAR>20 ⟹ simple model never called) \\
\texttt{gas\_exchange.oxygen\_concentration} & 210 & 210 & fixed & physical constant & Fixed atmospheric O2; excluded as non-varying \\
\texttt{carbohydrates.starch} & 0.05 & 1 & linear & initial condition & Observed range in spinach seedlings (\% DW; Geiger \& Servaites 1994) \\
\bottomrule
\end{tabular}
\end{table*}
